\subsection{Chronological surface and forecast evaluation}\label{sec:chronological_methods}

We extended the temporal evaluation in two stages to distinguish surface transfer from forecasting complete contract paths. The first stage used the completed expanding-window sector-network experiment on 58 capture dates through July 17, 2026. For each of 53 folds, the network fitted the first $k$ dates and predicted the next captured date, with $k=5,\ldots,57$. This experiment retained the original filters and capture-date labels, used two surface inputs, and calculated feature standardization from each training fold alone. Its 986{,}492 accepted observations contained no crossed or missing asks and no missing or nonpositive displayed bid or ask sizes under an independent reconstruction of the saved row counts. Since prediction used the later observation's stock price to calculate moneyness, the experiment measured conditional transfer of the surface across dates. It did not evaluate stock forecasts or the stochastic IV factor. The second stage fitted six sector networks once using 1{,}173{,}538 observations from 68 underlying sessions ending July 31. Their architectures and training schedule matched the existing two-input sector model, with seed 42 and a maximum of 2{,}000 epochs. Training-loss checkpoints, ticker levels, and training-only standardization were saved with the fitted models. The later sample contained 425{,}061 filtered IV observations across 23 sessions from August 4 through September 4. These observations supplied the final evaluation and did not select architectures, fitting windows, parameters, or seeds.

The second stage used an explicit capture audit before fitting or selecting contracts. The collector recorded retrieval times in UTC and stored the underlying-session close with each option row. We accepted retrievals after 20:00 UTC on that session and before 13:30 UTC on the next exchange session. This rule admitted overnight retrievals of a previous close while excluding twenty intraday repeat captures. We retained the latest eligible file per ticker and underlying session and one row per exact option symbol. Two-sided quotes required finite positive bids, noncrossed asks, and positive displayed sizes. Surface fitting additionally required the original IV and moneyness ranges and positive DTE, recomputed from expiration and the underlying session. Only files retrieved by the end of July 31 could enter training. Retrieval times were not exchange quote timestamps, so this audit did not establish quote freshness or exact stock--option synchronization. This limitation applied to the interpretation of both IV and dollar-price errors.

The forward test used GS and LLY at successive observed origins in August and early September, with planned horizons of one and five exchange sessions. Each origin supplied one put and one call with 25--45 calendar DTE. For each option type, we selected the expiration nearest 31 DTE and then the strike nearest $K/S=0.95$ for puts or $1.05$ for calls, breaking ties by earlier expiration and symbol. Selection used only the origin observations. We matched each selected symbol at its exact endpoint session and recorded missing endpoints without substituting another contract or capture date. The collection's rolling maturity buckets left none of these roughly monthly contracts observed five sessions later. We retained their one-session evaluation and added a separate shorter-maturity cohort with 10--21 DTE, choosing the expiration nearest 14 DTE with the same strike and tie rules. This availability-driven amendment preceded inspection of aggregate forecast errors and was recorded in the experiment protocol. The shorter cohort supplied the five-session case study; it did not replace the original contracts or validate their longer trajectories. All five IV variants began at the contract's observed origin IV. To compare their evolution without confounding the result with the fitted surface's initial level error, we rescaled the frozen variance target once at the origin:
\begin{equation}\label{eq:origin_anchor}
\bar v^{*}_{a,t}=\sigma_{a,0,\mathrm{obs}}^2
\frac{\bar v_a(S_t,\tau_t)}{\bar v_a(S_0,\tau_0)}.
\end{equation}
This anchoring defined a separate forecast experiment; the original sensitivity ablation retained its fitted initialization. Frozen IV held the common origin value, direct evaluation followed $\bar v^{*}_{a,t}$, and the three evolving factors used that anchored target with the previously specified parameters. Each origin used 1{,}000 JumpHMM paths and shared normal innovations across IV variants. The initial return state was drawn from the saved model's stationary distribution; recent returns did not update that state in the original forecast. An independent 10{,}000-path pilot per ticker supplied the growth shift to the existing 10\% annual prior and the one-session return normalization. Evaluation paths were not recentered. A zero-log-drift Gaussian random walk with the pilot return standard deviation provided a stock benchmark. Origin stock prices and IV updated as observations became available, but fitted parameters remained fixed. This was a retrospective evaluation at successive forecast origins, rather than a single forecast issued on July 31.

We scored forecast mean option values against the later bid--ask midpoints, using mean absolute error as the prespecified primary point metric. We retained that calculation and added median-based absolute errors to distinguish the choice of point summary from the IV comparison. Signed error, root mean squared error, and the fraction of predicted means inside the quoted spread supplied complementary price diagnostics. Distributional evaluation used the empirical continuous ranked probability score (CRPS)~\cite{gneiting2007}, which rewards concentration near the observed value while penalizing misplaced probability mass. For $N$ simulated values $P_j$ and observed midpoint $y$, we calculated the score as:
\begin{equation}\label{eq:forecast_crps}
\mathrm{CRPS}=\frac{1}{N}\sum_{j=1}^{N}|P_j-y|
-\frac{1}{2N^2}\sum_{j=1}^{N}\sum_{k=1}^{N}|P_j-P_k|.
\end{equation}
Smaller values indicate better distributional forecasts. We also reported central 90\% interval coverage and width. Scores were averaged across matched contracts within each ticker and origin before averaging dates, and stock forecast scores were reported separately. A diagnostic repeated IV evolution using the observed stock path only when every intervening session was available. It was explicitly conditional on realized stock information and therefore did not constitute a joint forecast. Repricing with the endpoint's reported IV in the same lattice supplied a further diagnostic of pricing discrepancies. All monetary scores were dollars per share. Overlapping horizons and common market movements made forecast dates and contracts dependent, so simulated path counts were not treated as independent market observations. The illustrative forecast was the earliest GS origin in the shorter-maturity cohort with both selected contracts observed at the five-session endpoint, regardless of its forecast errors.
